\section{RT-2 的结果：看见 novelty，也看见边界}

结果部分有四种证据：成功案例、分类汇总、ablation 和专门 benchmark。高质量读法是分别问它们回答什么，而不是把所有漂亮 demo 合并成“emergence”。本章按 qualitative、quantitative、causal 和 compositional 四层拆解。

\subsection{Emergent skills：组合新指令而非凭空创造}

Qualitative grid 展示 RT-2 对未在 robot data 中直接出现的对象、符号和语义组合执行动作，例如识别可乐品牌、选择最小物体或理解抽象类别。所谓 emergent 更接近\textbf{web semantics + learned motor primitive 的新组合}；动作空间仍来自训练过的 robot controller 与 token vocabulary。

\lecturefigure{slide-30-rt2-emergent-skills.jpg}{RT-2 emergent skills：符号、常识与新指令组合}{官方录像恢复幻灯片 V030；00:49:30}

\lecturefigure{slide-31-rt2-emergent-skills-chart.jpg}{Emergent skill 分类：symbol understanding、reasoning、human recognition}{官方录像恢复幻灯片 V031；00:50:24}

\begin{knowledgebox}{读图：平均柱状图会隐藏什么}
柱状图汇总多个 task family，先比较同一 category 内的模型，再看 average。若 reasoning 提升而 motor execution 不变，说明 semantic prior 在起作用；若 human recognition 高但 unseen environment 低，说明视觉概念迁移不等于场景鲁棒性。平均值不能告诉我们失败集中在哪些对象、动作或环境。
\end{knowledgebox}

\subsection{Seen / unseen 分类与 ablation}

Quantitative slide 分开 unseen objects、unseen backgrounds 和 unseen environments，这是必要设计：对象变化测试 semantic recognition，背景变化测试视觉鲁棒性，环境变化还会改变 geometry 与 affordance。Ablation slide 再比较 co-fine-tuning、只用 robot data、不同 model scale，尝试识别 transfer 来自哪里。

\lecturefigure{slide-32-rt2-quantitative-generalization.jpg}{RT-2 在 seen 与多类 unseen 设置中的成功率}{官方录像恢复幻灯片 V032；00:51:00}

\lecturefigure{slide-33-rt2-ablation-results.jpg}{RT-2 ablation：co-fine-tuning 与 data mixture 的贡献}{官方录像恢复幻灯片 V033；00:51:48}

\begin{warningbox}{Ablation 仍有外推限制}
若 ablation 在同一 robot、camera、task template 和 evaluation protocol 上进行，它能说明训练配置对该体系的影响，却不能直接回答新硬件、触觉控制、动态人类环境或数小时任务。实验控制越严格，因果解释越清楚；同时，适用范围也必须写清。
\end{warningbox}

\subsection{Language Table 与 chain-of-thought}

Language Table benchmark 测试由语言描述的桌面 manipulation，适合比较 instruction following 与组合泛化。Chain-of-thought 页面则让模型先生成视觉语义推理，再生成 action；它展示 intermediate reasoning 可以帮助处理“把能装水的物体移到某处”一类任务，但文字 reasoning 是否真实因果地驱动动作，仍需 intervention 而非只看自然语言 trace。

\lecturefigure{slide-34-language-table-benchmark.jpg}{Language Table：push-object instruction 的泛化结果}{官方录像恢复幻灯片 V034；00:52:03}

\lecturefigure{slide-35-rt2-chain-of-thought.jpg}{RT-2-PaLM-E 的 chain-of-thought 与行动序列}{官方录像恢复幻灯片 V035；00:53:30}

\begin{knowledgebox}{什么叫“reasoning 帮助了 policy”}
最强证据不是模型写出了看似合理的解释，而是改变或移除 intermediate reasoning 会系统性改变 action success；或者 model 在需要多步关系组合的 held-out tasks 上显著优于同容量、不生成 reasoning 的 baseline。自然语言 trace 是诊断窗口，不自动等于 faithful internal computation。
\end{knowledgebox}

\subsection{VLA 总结与仍未解决的问题}

VLA summary 页把新任务、新对象、chain-of-thought 与 generalization 列为收益，同时把 data diversity、额外 context 与 capability 扩展列为未来方向。它没有宣告 low-level control 已解决；恰恰相反，结果表明 semantic transfer 已出现，但 robot data、动作精度和 closed-loop reliability 仍是决定性瓶颈。

\lecturefigure{slide-36-vla-summary.jpg}{VLA 模型的能力总结与未来问题}{官方录像恢复幻灯片 V036；00:55:36}

\teachervoice{讲者把这些结果称为“semantic transfer”而不是通用机器人智能。课堂语气很重要：模型能执行训练中没有逐字出现的命令，是值得重视的新能力；但它仍在相同 action space、机器人平台和有限场景中工作。}

\subsection{本章小结}

RT-2 的可信结论来自多层证据共同支持：qualitative cases 展示新组合，category metrics 衡量泛化，ablation 追踪 co-fine-tuning，Language Table 测 instruction generalization。它们支持 VLM prior 可迁移到控制，但没有消除 embodiment、可靠性与安全边界。

